\section{Conclusión}

En esta práctica vimos cómo funciona IPsec y cómo se usa para armar una VPN
entre dos redes que están conectadas por un internet que no es seguro. Primero
comprobamos que, sin VPN, el sniffer ve los pings (ICMP) tal cual, con
direcciones y todo, lo que muestra lo fácil que es espiar una comunicación que
no va protegida.

Después configuramos la VPN en los routers R\_01 y R\_03 y entendimos que el
proceso tiene dos fases: en la Fase 1 se crea un canal seguro para ponerse de
acuerdo (ISAKMP), y en la Fase 2 se crea el túnel por donde realmente viajan los
datos cifrados (IPsec). También vimos que la configuración tiene que coincidir en
los dos extremos (cifrado, hash, grupo Diffie-Hellman y llave compartida) o el
túnel no se levanta. Ya con la VPN activa, el sniffer solo mostró paquetes ESP,
así que el contenido y las IPs privadas de las laptops quedaron ocultos.

Con las preguntas aprendimos cosas importantes: que conviene cambiar las llaves
seguido para no darle tiempo a un atacante, que los grupos Diffie-Hellman viejos
(2 y 5) ya no son seguros y hay que pasarse al 14 o mejor a curvas elípticas, y
que Perfect Forward Secrecy protege lo que ya se mandó aunque en el futuro
roben una llave.

En resumen, una VPN con IPsec nos da confidencialidad, integridad y autenticidad,
y sirve para conectar sedes sin que nadie en el camino pueda ver lo que se
manda. Pero no basta con activarla: hay que elegir bien los parámetros para que
de verdad sea segura.
